\chapter{Cross-Sectional Dependence}
\label{ch:cross}

The serial-dependence layer of Chapter~\ref{ch:serial} delivers the pairwise standard
errors $\seNW{jk}$, but the simulation-based critical value needs a full covariance of
the contrasts. I build that covariance in two ways: by reconstructing a level covariance
from the pairwise standard errors (Approach~1, classical multidimensional scaling), and
by estimating the covariance of the differences directly (Approach~2, the object the
procedure uses). I prove that the two coincide on a balanced panel, exactly in finite
samples under a common bandwidth and asymptotically otherwise. On an unbalanced panel
they target different objects, and only the direct construction is the covariance the
procedure needs.

Figure~\ref{fig:contrast-covariance-routes} traces both constructions from the pairwise
loss differences to the common inference procedure of Chapter~\ref{ch:downstream}.

\figContrastCovarianceRoutes

\section{The calibration condition}
\label{sec:calibration}

The critical value comes from simulating the joint distribution of the studentized
contrasts under the null $\Delta_{j,k}=0$. The simulation draws a Gaussian vector
$Z\sim\Normal(0,\Sh)$ and forms $(Z_j - Z_k)/\seNW{jk}$ for each pair. For these draws
to share the distribution of the real studentized statistics
$T_{jk} = (\hat{\Delta}_{j,k}-\Delta_{j,k})/\seNW{jk}$, the variance of the simulated
numerator must equal the squared denominator:
\begin{equation}
  \Var(Z_j - Z_k)
  \;=\; \Sh_{jj} + \Sh_{kk} - 2\Sh_{jk}
  \;\stackrel{!}{=}\; \big(\seNW{jk}\big)^2
  \qquad \forall\, j\neq k.
  \label{eq:calib}
\end{equation}
So I need a covariance matrix $\Sh$ that satisfies \eqref{eq:calib}. The naive approach,
estimating $\Sh$ from level variances, fails on an unbalanced panel. The level means are
taken over each forecaster's own dates, so their differences are not the overlap
contrasts $\hat\Delta_{j,k}$, and the common-shock term $c_t$ of \eqref{eq:decomp} no
longer cancels in them. Rather than estimate $\Sh$ and hope that \eqref{eq:calib} holds,
Approach~1 imposes \eqref{eq:calib} and backs out $\Sh$; Approach~2 sidesteps the level
covariance entirely.

\section{Approach 1: the MDS reconstruction}
\label{sec:mds}

The left-hand side of \eqref{eq:calib} has the form of a squared Euclidean distance. If
$\Sh$ is an inner-product (Gram) matrix, i.e.\ $\Sh_{jk} = \langle v_j,
v_k\rangle$ for some configuration $v_1,\dots,v_p$, then
\begin{equation}
  \|v_j - v_k\|^2
  = \langle v_j,v_j\rangle + \langle v_k,v_k\rangle - 2\langle v_j,v_k\rangle
  = \Sh_{jj} + \Sh_{kk} - 2\Sh_{jk}.
\end{equation}
A Gram matrix $G = V V^\top$ is positive semi-definite and encodes an inner-product
space. So \eqref{eq:calib} asks for a Gram matrix whose induced squared distances equal
$S^2_{jk} = (\seNW{jk})^2$. This is the \emph{classical multidimensional scaling}
problem: recover a point configuration from its pairwise distances. It has a
closed-form solution (Torgerson's classical scaling), and Schoenberg's theorem
characterises when that solution is valid.

\subsection{The double-centring formula}

Let $S^2$ be the $(p\times p)$ matrix of squared difference standard errors (zero
diagonal), and let $J \equiv I_p - \tfrac{1}{p}\,\one\one^\top$ be the centring matrix.
The construction is
\begin{equation}
  \boxed{\;\Sh \;=\; -\tfrac{1}{2}\, J\, S^2\, J.\;}
  \label{eq:doublecentre}
\end{equation}

\paragraph{Calibration holds by construction.}
Double-centring acts entrywise as
\begin{equation}
  \Sh_{jk}
  = -\tfrac12\big(S^2_{jk} - \bar{s}_{j\cdot} - \bar{s}_{\cdot k}
      + \bar{s}_{\cdot\cdot}\big),
\end{equation}
with row mean $\bar{s}_{j\cdot}$, column mean $\bar{s}_{\cdot k}$, and grand mean
$\bar{s}_{\cdot\cdot}$. Using $S^2_{jj}=0$ and symmetry,
$\Sh_{jj}=\bar{s}_{j\cdot}-\tfrac12\bar{s}_{\cdot\cdot}$ and
$\Sh_{kk}=\bar{s}_{k\cdot}-\tfrac12\bar{s}_{\cdot\cdot}$. In the contrast, the row,
column, and grand means all cancel:
\begin{align}
  \Sh_{jj} + \Sh_{kk} - 2\Sh_{jk}
  &= \Big(\bar{s}_{j\cdot} - \tfrac12\bar{s}_{\cdot\cdot}\Big)
   + \Big(\bar{s}_{k\cdot} - \tfrac12\bar{s}_{\cdot\cdot}\Big)
   + \big(S^2_{jk} - \bar{s}_{j\cdot} - \bar{s}_{\cdot k}
          + \bar{s}_{\cdot\cdot}\big) \notag\\[2pt]
  &= S^2_{jk}. \qquad\checkmark
\end{align}
The identity holds for any symmetric, zero-diagonal $S^2$, whether or not it is
Euclidean-embeddable. Calibration is therefore guaranteed, not estimated. Because only
the difference standard errors enter, the common shock never enters $\Sh$. The diagonal
$\Sh_{jj}$ is defined by the very contrasts the procedure tests, not by tail-dominated
marginal variances.

\subsection{Why $\Sh$ is rank-deficient (and why that is correct)}

Because $J\one = 0$, it follows that $\Sh\,\one = -\tfrac12\, J S^2 (J\one) = 0$, so the
constant vector $\one$ lies in the null space of $\Sh$: the matrix is singular by
construction. In probability terms, taking $w=\one$ in $\Var(w^\top Z)=w^\top\Sh\,w$
gives $\Var(Z_1+\dots+Z_p)=\one^\top\Sh\,\one=0$. The overall level is pinned to zero,
and $Z$ never moves along the all-ones direction. This is the statement that only
contrasts are identified. Because $\Sh$ is built from differences, every quantity the
method uses is blind to a common shift $Z_j \mapsto Z_j + c$. A full-rank $\Sh$ would be
wrong here: it would assign variance to a level the ranking can never observe.
Simulating from a singular PSD matrix is standard (draw on the non-zero
eigen-directions), and the procedure never inverts $\Sh$, so the singularity costs
nothing.

\subsection{The nearest-PSD projection}
\label{sec:mds-psd}

Double-centring always delivers calibration, but $\Sh$ is a valid covariance only if it
is positive semi-definite. By Schoenberg's theorem, $-\tfrac12 J S^2 J$ is PSD if and
only if $S^2$ is Euclidean-embeddable. The entries of $S^2$ are estimated pairwise, each
on its own overlap and bandwidth, so embeddability generally fails and $\Sh$ typically
has a few small negative eigenvalues. The repair is the nearest-PSD projection
\citep{higham1988}: eigendecompose $\Sh = Q\Lambda Q^\top$, clip
$\Lambda^{+} = \max(\Lambda, 0)$, and reassemble $\Sh^{+} = Q\Lambda^{+}Q^\top$. The
result is the closest PSD matrix to $\Sh$ in Frobenius norm. If $S^2$ is embeddable, the
clip does nothing.

\begin{remark}[Calibration vs.\ PSD-ness]
The clip perturbs calibration: the identity $\Sh_{jj}+\Sh_{kk}-2\Sh_{jk} = S^2_{jk}$
holds exactly only for the unprojected $\Sh$. Eigenvalue clipping trades a small
calibration error for validity. Because $\Sh^{+}-\Sh$ is positive semi-definite with
largest eigenvalue $|\lambda_{\min}(\Sh)|$, the clip can only inflate the simulated
contrast variances, by at most $2|\lambda_{\min}|$, which makes the critical value
conservative. The clipped eigenvalues grow with the imbalance of the panel and with the
heterogeneity of the bandwidths; on a mildly unbalanced panel with similar persistences
they are approximately zero. The Frobenius norm of the repair,
$\Frob{\Sh - \Sh^{+}}$, serves as a robustness diagnostic.
\end{remark}

\subsection{Non-uniqueness of the calibrated covariance}
\label{sec:nonuniqueness}

The double-centred matrix is not the only covariance that satisfies \eqref{eq:calib}.
Condition \eqref{eq:calib} imposes $\binom{p}{2}$ equations on the $\binom{p}{2}+p$ free
entries of a symmetric $\Sigma$, so its solution set is a $p$-dimensional affine family,
of which $-\tfrac12 J S^2 J$ is one canonical member. If $\Sigma$ satisfies
\eqref{eq:calib}, so does
\begin{equation}
  \Sigma' = \Sigma + a\one^\top + \one a^\top
\end{equation}
for any $a\in\Real^p$, since $a\one^\top + \one a^\top$ contributes
$2a_j+2a_k-2a_j-2a_k=0$ to every difference variance. Adding such a term reparametrises
the unobservable levels, to which all contrasts are invariant. Within this family the
double-centred solution is the unique \emph{mean-centred} representative
($\one^\top\Sh=0$), the gauge that places the recovered configuration at the origin.
The ambiguity is harmless. The simulation uses only the contrasts $\{Z_j-Z_k\}$, whose
joint distribution is the same across the family. And the mean-centred representative is
the right input to the projection, because it parks the unidentified level in the null
space instead of smearing it across the spectrum (a generic shift can drive eigenvalues
negative).

\section{Approach 2: the direct difference covariance $\OmegaDirect$}
\label{sec:omega}

Approach~1 reconstructs a $p$-dimensional level covariance and differences it.
Approach~2 models directly the object the procedure uses: the covariance of the
$q=\binom{p}{2}$ difference-mean estimators.

\subsection{The balanced case}
\label{sec:omega-balanced}

Stack the difference series into a $q$-vector process $d_t = (d^a_t)_{a=1}^{q}$, with
$a=(j,k)$, $j<k$, and $d^a_t = X_{t,j}-X_{t,k}$. When the panel is balanced, so that
every difference is observed on the same $n$ dates, a single multivariate Newey--West
estimator gives the covariance of the difference-mean vector
$\hat\Delta = \tfrac1n\sum_t d_t$ in one pass:
\begin{equation}
  \Omegahat \;=\; \hat S/n,
  \qquad
  \hat S = \Gammahat_0 + \sum_{h=1}^{L} w_h\big(\Gammahat_h + \Gammahat_h^\top\big),
  \qquad
  \Gammahat_h = \tfrac1n\sum_t \tilde d_t\,\tilde d_{t-h}^\top,
  \label{eq:omega-balanced}
\end{equation}
where $\tilde d_t = d_t - \bar d$ and $w_h$ is the Bartlett weight of
Chapter~\ref{ch:serial}. The estimator is positive semi-definite by construction
(Bartlett lies in the Andrews class $\mathcal{K}_2$). Calibration is automatic, since its
diagonal is the univariate long-run variance of each $d^a$. The common shock never
enters, since the estimator operates on differences. Its theoretical home is
\citet{andrews1991}, which is natively multivariate.

\subsection{The unbalanced case: the pairwise-of-pairs estimator}
\label{sec:omega-unbal}

The balanced construction requires every difference series to be observed on the same
dates. Under staggered entry and exit the common window $\bigcap_a T_a$ is small or
empty, and \eqref{eq:omega-balanced} is infeasible. I instead estimate each entry of
$\OmegaDirect$ on the dates its two difference series share (the pairwise-of-pairs
route). Write $d^a_t$ on its support $T_a = T_j\cap T_k$ with $n_a=|T_a|$. The entry is
the sampling covariance of two difference-mean estimators,
$\Omega_{ab}=\Cov(\hat\Delta_a,\hat\Delta_b)$, and the lag-$h$ cross-covariance is
identified only where both series are observed, on the paired index set
\begin{equation}
  T_{ab}(h) = \{\, t : t\in T_a,\ t-h\in T_b \,\},
  \qquad
  n_{ab} = |T_{ab}(0)| = |T_a\cap T_b|.
\end{equation}
The pair-of-pairs overlap $n_{ab}$ is thus the sample size on which the cross-covariance
of $d^a$ and $d^b$ is estimated. It is not the sample size that scales the entry.
$\hat\Delta_a$ averages over $n_a$ dates and $\hat\Delta_b$ over $n_b$, and only the
dates they share (and their near neighbours, under serial dependence) contribute to
their covariance:
\begin{equation}
  \Omega_{ab}
  = \frac{1}{n_a n_b}\sum_{t\in T_a}\sum_{s\in T_b}\Cov\big(d^a_t,d^b_s\big)
  \;\approx\; \frac{n_{ab}}{n_a\,n_b}\;\omega_{ab},
  \label{eq:omega-scaling}
\end{equation}
with $\omega_{ab}$ the long-run cross-covariance of $d^a$ and $d^b$. This is the familiar
covariance of two means computed on partially overlapping samples; it reduces to
$\omega_{ab}/n$ when $T_a=T_b$ has $n$ dates. With each series demeaned on its own
support, $\tilde d^a_t = d^a_t - \bar d^a$, the Bartlett-weighted HAC estimator of the
entry is
\begin{equation}
  \boxed{\;
  \Omegahat_{ab}
  = \frac{n_{ab}}{n_a\,n_b}\left[\hat\gamma^{ab}_0
    + \sum_{h=1}^{L_{ab}} w_h\big(\hat\gamma^{ab}_h + \hat\gamma^{ba}_h\big)
    \right],
  \qquad
  \hat\gamma^{ab}_h
  = \frac{1}{n_{ab}}\!\!\sum_{t\in T_{ab}(h)}\!\! \tilde d^a_t\,\tilde d^b_{t-h},
  \qquad
  w_h = 1-\frac{h}{L_{ab}+1}.
  \;}
  \label{eq:omega-entry}
\end{equation}
Both lag directions appear because cross-covariances are not lag-symmetric
($\gamma^{ab}_{-h}=\gamma^{ba}_h$). Two design choices settle the estimator.

\paragraph{Bandwidth.} Each difference series has its own persistence, so I set the
bandwidth per series, $L_a$, by the Andrews plug-in of \S\ref{sec:bandwidth}. An
off-diagonal entry needs one bandwidth for the shared lag sum; I take
$L_{ab}=\min\{L_a,L_b\}$, the more cautious of the two choices on short overlaps. The
rule reduces to $L_{aa}=L_a$ on the diagonal. The implementation further caps $L_{ab}$ at
$n_{ab}-1$ and sets to zero any off-diagonal entry whose overlap $n_{ab}$ is below a
minimum (two dates). The rule has a cost. The lagged cross-covariances of $d^a$ and $d^b$
are not bounded by either series' own persistence, so if one of them is serially
uncorrelated (and its plug-in bandwidth stays bounded) while the other is persistent,
the truncated cross sum need not converge to the long-run cross-covariance. The diagonal
entries are unaffected.

\paragraph{Normalization.} Multiplying out, $\Omegahat_{ab} = s_{ab}/(n_a n_b)$, where
$s_{ab}$ is the Bartlett-weighted cross sum over jointly observed dates. Written this
way, the entry is defined, and zero, when the two supports never meet. The divisor is a
single number per entry, applied uniformly across all lags rather than by the per-lag
count $|T_{ab}(h)|$, and it is the product of the two support sizes rather than the
square of their overlap. Both choices matter for positive semi-definiteness. Fill each
demeaned series with zeros off its support and stack them into $y_t\in\Real^q$. The
cross sums over jointly observed dates are then exactly the full-sample sums of the
zero-filled series, so under a common bandwidth $L$ the matrix
$S=(s_{ab})$ is the Bartlett quadratic form
$\sum_t y_t y_t^\top + \sum_{h=1}^{L} w_h \sum_t (y_t y_{t-h}^\top + y_{t-h}y_t^\top)$,
which is positive semi-definite, and
\begin{equation}
  \Omegahat = N^{-1} S\, N^{-1},
  \qquad N = \operatorname{diag}\big(n_a\big)_{a=1}^{q},
  \label{eq:omega-congruence}
\end{equation}
is a congruence of it, hence also positive semi-definite on any unbalanced panel,
provided no entry is zeroed or truncated by the implementation rules above, that is,
provided every pair of pairs whose supports come within $L$ periods of each other shares
more than $\max(L,1)$ dates. A per-lag normalization scales each lag block differently
and destroys the quadratic form. Dividing by $n_{ab}^2$, the natural-looking choice, is
not a congruence either. It also overstates every off-diagonal entry by the factor
$n_a n_b/n_{ab}^2\ge 1$, so the implied correlations between contrasts are too strong
whenever supports differ.

\subsection{Nesting and diagonal-consistency}
\label{sec:omega-nesting}

With this normalization the estimator ties back to the rest of the pipeline.
\emph{Nesting}: when the panel is balanced and a common bandwidth $L$ is used,
$n_a=n_{ab}=n$ for every $a,b$, and each entry \eqref{eq:omega-entry} coincides with the
$(a,b)$ entry of the stacked estimator \eqref{eq:omega-balanced}. The pairwise-of-pairs
route is therefore a strict generalization, not a different method.
\emph{Diagonal-consistency}: on the diagonal $a=b$, \eqref{eq:omega-entry} reduces to
the univariate Bartlett long-run variance of $d^a_t$ divided by $n_a$. When the support
$T_a$ is a contiguous block of periods, as on a balanced panel or under staggered entry
and exit, this is exactly the squared serial-dependence standard error,
$\Omegahat_{aa} = (\seNW{a})^2$. When $T_a$ has gaps, the two differ slightly.
\eqref{eq:omega-entry} pairs periods that are $h$ calendar periods apart, whereas
$\seNW{a}$ counts lags in common observations (\S\ref{sec:nw}), so autocovariance terms
that straddle a gap enter one but not the other (Appendix~\ref{app:impl-details}
quantifies the difference on the NGDP panel). Either way, calibration of the direct
construction is automatic, because it studentizes by its own diagonal. The simulation
denominator $\sqrt{\Omegahat_{aa}}$ and the variance of the simulated numerator are the
same diagonal entry; after the projection of \S\ref{sec:omega-psd}, they are the same
entry of the projected matrix.

\subsection{Nearest-PSD projection}
\label{sec:omega-psd}

By \eqref{eq:omega-congruence}, $\Omegahat$ is positive semi-definite under a common
bandwidth, balanced or not. The pipeline, however, uses the per-series Andrews bandwidth
of \S\ref{sec:bandwidth} and, off the diagonal, the shorter of the two,
$L_{ab}=\min\{L_a,L_b\}$. Entries assembled at different bandwidths no longer form a
single quadratic form. On a balanced panel the matrix is then generically indefinite,
with small negative eigenvalues, because its target has rank $p-1<q$ and sits on the
boundary of the positive semi-definite cone. On an unbalanced panel it can be indefinite
as well (\S\ref{sec:caveat}). The repair is again the Frobenius-nearest PSD projection
\citep{higham1988}, applied directly to the modelled object $\Omegahat$. I report
$\Frob{\Omegahat-\Omegahat^{+}}$ as a diagnostic of how far the heterogeneous bandwidths
push the estimate from positive semi-definiteness.

\section{Equivalence of the two constructions}
\label{sec:equivalence}

Let $A$ be the fixed $q\times p$ contrast matrix whose row $a=(j,k)$ has $+1$ in column
$j$, $-1$ in column $k$, and zero elsewhere, so that $D = A Z$ is the vector of pairwise
differences of a level vector $Z$. Approach~1 draws $Z\sim\Normal(0,\Sh^{+})$ and forms
$D=AZ\sim\Normal(0,A\Sh^{+} A^\top)$; Approach~2 draws $D\sim\Normal(0,\Omegahat^{+})$
directly. The two are the same test when $A\Sh A^\top$ and $\Omegahat$ agree, together
with the standard errors that studentize the draws. Everything rests on one algebraic
fact: a covariance of the form $A\Sigma A^\top$ is determined by its diagonal.

\begin{lemma}[Polarization]
\label{lem:polarization}
Let $\Sigma$ be any symmetric $p\times p$ matrix and
$S^2_{jk}=\Sigma_{jj}+\Sigma_{kk}-2\Sigma_{jk}$. For pairs $a=(j,k)$ and $b=(l,m)$,
\begin{equation}
  \big(A\Sigma A^\top\big)_{ab}
  = \Sigma_{jl}-\Sigma_{jm}-\Sigma_{kl}+\Sigma_{km}
  = \tfrac12\big(S^2_{jm}+S^2_{kl}-S^2_{jl}-S^2_{km}\big).
  \label{eq:polarization}
\end{equation}
Moreover $-\tfrac12 J S^2 J = J\Sigma J$, and therefore
$A\big(-\tfrac12 J S^2 J\big)A^\top = A\Sigma A^\top$.
\end{lemma}

\begin{proof}
The first equality is the covariance of the contrasts $Z_j-Z_k$ and $Z_l-Z_m$ when
$\Var(Z)=\Sigma$. For the second, substitute the definition of $S^2$. The diagonal terms
cancel in pairs,
$(\Sigma_{jj}+\Sigma_{mm})+(\Sigma_{kk}+\Sigma_{ll})-(\Sigma_{jj}+\Sigma_{ll})
-(\Sigma_{kk}+\Sigma_{mm})=0$, leaving
$-2\Sigma_{jm}-2\Sigma_{kl}+2\Sigma_{jl}+2\Sigma_{km}$; the identity also holds when
the pairs share an index, with $S^2_{jj}=0$. For the last statement, write
$S^2=\sigma\one^\top+\one\sigma^\top-2\Sigma$ with $\sigma=(\Sigma_{11},\dots,
\Sigma_{pp})^\top$. Since $J\one=0$, the first two terms vanish under double-centring,
so $-\tfrac12 JS^2J=J\Sigma J$. Finally, $A\one=0$ gives $AJ=A$, hence
$AJ\Sigma JA^\top=A\Sigma A^\top$.
\end{proof}

The lemma makes precise what the MDS route recovers. It does not recover $\Sigma$: the
level of the losses, and with it the common shock, is not identified from contrasts. It
recovers only $J\Sigma J$, and the contrast matrix cannot tell the two apart. Given the
pairwise variances of any matrix of the form $A\Sigma A^\top$, the MDS route reproduces
that matrix entry by entry. Whether it reproduces the covariance the procedure needs
therefore turns on whether that covariance has the form $A\Sigma A^\top$ at all.

\begin{proposition}[Equivalence on a balanced panel]
\label{prop:equivalence}
Suppose every forecaster is observed on the same $n$ dates, so that
$\hat\Delta=A\hat\theta$ with $\hat\theta$ the vector of mean losses, and let
$\Sigma_n=\Var(\hat\theta)$.
\begin{enumerate}[label=(\roman*)]
  \item \emph{Population.} $\OmegaDirect=\Var(\hat\Delta)=A\Sigma_nA^\top$, and the MDS
    map applied to the true pairwise variances returns it:
    $A\big(-\tfrac12 JS^2J\big)A^\top=\OmegaDirect$.
  \item \emph{Finite samples, common bandwidth.} If every pairwise standard error and
    $\Omegahat$ use the same Bartlett bandwidth $L$, then $\Sh=-\tfrac12 J\hat S^2J$ is
    positive semi-definite, neither construction is projected, $A\Sh A^\top=\Omegahat$,
    and $\seNW{a}=\sqrt{\Omegahat_{aa}}$ for every $a$. The two constructions produce
    identically distributed studentized draws: they are the same test.
  \item \emph{Asymptotically, per-series bandwidths.} If $n\Sigma_n\to\Sigma_\infty$ with
    $A\Sigma_\infty A^\top$ having a positive diagonal, and every pairwise standard error
    and every entry of $\Omegahat$ is consistent after scaling by $n$, i.e.\
    $n(\seNW{a})^2\to_p(A\Sigma_\infty A^\top)_{aa}$ and
    $n\Omegahat_{ab}\to_p(A\Sigma_\infty A^\top)_{ab}$ --- which the Andrews plug-in
    delivers under the conditions of \S\ref{sec:bandwidth} when every difference series
    is persistent, so that every $L_{ab}\to\infty$; if some $L_a$ stays bounded, the
    off-diagonal entries involving $a$ are consistent only if the corresponding
    cross-covariances vanish at the truncated lags --- then
    $nA\Sh^{+}A^\top$ and $n\Omegahat^{+}$ converge in probability to the same limit
    $A\Sigma_\infty A^\top$, and the critical values of the two constructions converge to
    the same limit.
\end{enumerate}
\end{proposition}

\begin{proof}
(i) On a balanced panel $\hat\Delta_a=\hat\theta_j-\hat\theta_k$ for every pair, so
$\hat\Delta=A\hat\theta$ and $\OmegaDirect=A\Sigma_nA^\top$. The true pairwise variances
are its diagonal, $S^2_{jk}=(\Sigma_n)_{jj}+(\Sigma_n)_{kk}-2(\Sigma_n)_{jk}$, and
Lemma~\ref{lem:polarization} with $\Sigma=\Sigma_n$ gives the claim.

(ii) Let $x_t\in\Real^p$ be the losses at date $t$ and $\tilde x_t=x_t-\bar x$. The
difference vector is $d_t=Ax_t$, and demeaning commutes with $A$, so
$\tilde d_t=A\tilde x_t$ and every sample autocovariance factors,
$\Gammahat_h(d)=A\Gammahat_h(x)A^\top$. On a balanced panel $\Omegahat$ reduces to the
stacked estimator \eqref{eq:omega-balanced} (nesting, \S\ref{sec:omega-nesting}), hence
$\Omegahat=A\hat VA^\top$ with $\hat V$ the $p$-variate Bartlett Newey--West estimator of
$\Var(\hat\theta)$, which is positive semi-definite. By diagonal-consistency (the
supports are contiguous) the diagonal of $\Omegahat$ is $(\seNW{a})^2$, so
$\hat S^2_{jk}=\hat V_{jj}+\hat V_{kk}-2\hat V_{jk}$. Lemma~\ref{lem:polarization} with
$\Sigma=\hat V$ gives $\Sh=J\hat VJ$, which is positive semi-definite as a congruence of
$\hat V$, so its projection is the identity. And $A\Sh A^\top=A\hat VA^\top=\Omegahat$,
which is positive semi-definite and likewise unprojected. The MDS draws
$(Z_j-Z_k)/\seNW{a}$ with $Z\sim\Normal(0,\Sh)$ are $AZ\sim\Normal(0,\Omegahat)$
divided by $\seNW{a}=\sqrt{\Omegahat_{aa}}$, which is the direct construction's
$D_a/\sqrt{\Omegahat_{aa}}$ in distribution. The observed statistics share the same
denominators.

(iii) By consistency, $n\hat S^2$ converges to the matrix of squared contrast distances
of $\Sigma_\infty$. Double-centring is linear and the nearest-PSD projection is
continuous and positively homogeneous, so $nA\Sh^{+}A^\top=A\big(-\tfrac12 J(n\hat S^2)J
\big)^{+}A^\top$ converges to $A(J\Sigma_\infty J)^{+}A^\top$. The limit $J\Sigma_\infty
J$ is positive semi-definite, so the projection acts as the identity there, and by
Lemma~\ref{lem:polarization} the limit is $A\Sigma_\infty A^\top$. Likewise
$n\Omegahat\to A\Sigma_\infty A^\top$, a positive semi-definite limit, so
$n\Omegahat^{+}$ has the same limit. Both studentizers, $n(\seNW{a})^2$ and
$n\Omegahat^{+}_{aa}$, converge to the positive $(A\Sigma_\infty A^\top)_{aa}$. Both
vectors of studentized draws therefore converge in distribution to $\Normal(0,R)$, with
$R$ the correlation matrix of $A\Sigma_\infty A^\top$. The maximum of the signed entries
of a centred Gaussian vector with positive-variance components, over any non-empty set
of ordered contrasts, has a continuous distribution function that is strictly increasing
on $(0,\infty)$. Hence the critical values for any given set of active contrasts converge
to the same limit.
\end{proof}

\begin{remark}[Unbalanced panels: two different targets]
\label{rem:unbalanced-targets}
Proposition~\ref{prop:equivalence} uses balance at its first step. When the supports
differ, $\hat\Delta_a$ is the mean of $d^a$ over its own support $T_a$, and in general
$\hat\Delta$ is not $A\hat\theta$ for any vector $\hat\theta$. The combination
$\hat\Delta_{12}+\hat\Delta_{23}-\hat\Delta_{13}$, identically zero on a balanced panel,
is now a non-degenerate combination of means over three different sets of dates. The
target $\OmegaDirect=\Var(\hat\Delta)$ then (for $p\ge3$) generally has full rank
$q>p-1$, whereas every matrix $A\Sigma A^\top$ has rank at most $p-1$. No level
covariance reproduces $\OmegaDirect$, and \eqref{eq:polarization} fails for it: its
off-diagonal entries are not functions of its diagonal. The MDS route nevertheless sets
every off-diagonal entry from the pairwise variances through \eqref{eq:polarization}. On
an unbalanced panel it therefore estimates a different object from $\OmegaDirect$ even
with consistent inputs. Moreover, the true pairwise variances need not be
Euclidean-embeddable, and that failure is part of what its projection repairs. The
direct estimator estimates each entry from the data and, by \eqref{eq:omega-scaling},
targets $\OmegaDirect$ itself (subject to the bandwidth caveat of
\S\ref{sec:omega-unbal}).
\end{remark}

On a balanced panel, then, the two constructions are the same test, and Approach~2 is
the cleaner route to it: it is estimable directly from the data, its calibration is
automatic rather than imposed, and it avoids the double-centring reconstruction and the
non-uniqueness of \S\ref{sec:nonuniqueness}. On an unbalanced panel Approach~2 is not
merely the better-calibrated of two routes to a common answer. It is the only one of the
two that targets the covariance of the contrasts the procedure tests.

\section{What the direct approach buys, and when it needs the projection}
\label{sec:caveat}

Three properties of the direct construction are easily conflated. The first is
\emph{positive semi-definiteness}. By \eqref{eq:omega-congruence}, $\Omegahat$ is
positive semi-definite under a common bandwidth on balanced and unbalanced panels alike;
imbalance by itself does not break the guarantee. Heterogeneity of the bandwidth does.
With the per-series Andrews bandwidth of \S\ref{sec:bandwidth}, entries assembled at
different bandwidths $L_{ab}=\min\{L_a,L_b\}$ no longer form a single quadratic form, and
the matrix is generically indefinite even on a balanced panel; its negative eigenvalues
remain small when the persistences of the difference series are similar. The guarantee
should therefore be stated as ``common bandwidth'', not ``balanced panel''.

The second is \emph{calibration}. The direct construction studentizes by its own
diagonal, so its draws are unit-variance by construction, whereas the MDS route imposes
calibration before its projection and loses it to the extent that the projection
intervenes (\S\ref{sec:mds-psd}). The third is the \emph{target}: on an unbalanced
panel only the direct construction estimates the covariance of the contrasts
(Remark~\ref{rem:unbalanced-targets}). Both constructions use the nearest-PSD
projection, but on different objects. Approach~2 applies it to the modelled object
$\Omegahat$, where it repairs only the bandwidth heterogeneity (including the truncation
$L_{ab}\le n_{ab}-1$ and the zeroed thin overlaps). Approach~1 applies it to a
reconstructed level covariance, where on an unbalanced panel it must also absorb the
mismatch between the pairwise variances and any configuration of points.
